\section*{第 \ref{ch:g8:lenses-images} 章——透镜与像}

\begin{solution}{exo:g8:lenses-images:1}
中间鼓起：\omterm{def:g8:lenses-images:families}{凸透镜}——它使平行\omterm{def:g6:light-propagation:ray}{光束}的\omterm{def:g6:light-propagation:ray}{光线}相互靠拢，
会聚于\omterm{def:g8:lenses-images:focus}{焦点}。中间凹陷：\omterm{def:g8:lenses-images:families}{凹透镜}——
它使\omterm{def:g6:light-propagation:ray}{光束}向外散开。
\end{solution}

\begin{solution}{exo:g8:lenses-images:2}
\omterm{def:g8:lenses-images:focus}{焦点}是平行\omterm{def:g6:light-propagation:ray}{光束}经过\omterm{def:g8:lenses-images:families}{透镜}后会聚的点，\omterm{def:g8:lenses-images:focus}{焦距} $f$ 即
\omterm{def:g8:lenses-images:families}{透镜}到\omterm{def:g8:lenses-images:focus}{焦点}的距离。$f$ 越小，\omterm{def:g8:lenses-images:families}{透镜}弯曲越厉害，
对\omterm{def:g6:light-propagation:ray}{光线}的偏折作用越强。
\end{solution}

\begin{solution}{exo:g8:lenses-images:3}
把远处景物的像投射到墙上，移动\omterm{def:g8:lenses-images:families}{透镜}使像清晰，
再量取\omterm{def:g8:lenses-images:families}{透镜}到墙的距离，这个距离就是 $f$。景物
必须足够远，使其\omterm{def:g6:light-propagation:ray}{光线}近似平行——
平行\omterm{def:g6:light-propagation:ray}{光线}按\omterm{def:g8:lenses-images:focus}{焦点}的定义，恰好会聚于
距\omterm{def:g8:lenses-images:families}{透镜} $f$ 处。
\end{solution}

\begin{solution}{exo:g8:lenses-images:4}
切勿把阳光聚焦到易燃物附近——\omterm{def:g8:lenses-images:focus}{焦点}处的光斑
几秒钟内就会烧焦并引燃。切勿透过任何\omterm{def:g8:lenses-images:families}{透镜}直视太阳——
眼睛加上\omterm{def:g8:lenses-images:families}{透镜}，会把灼热的光斑聚焦到视网膜上，
瞬间造成伤害，而且毫无痛感。
\end{solution}

\begin{solution}{exo:g8:lenses-images:5}
由会聚\omterm{def:g6:light-propagation:ray}{光线}真正形成的像，可以用屏幕承接。
其特征是倒立且明亮——\omterm{def:g8:lenses-images:families}{透镜}在每一点
收集的光可达整个\omterm{def:g8:lenses-images:families}{透镜}的量，而小孔
只让一束\omterm{def:g6:light-propagation:ray}{光线}通过。
\end{solution}

\begin{solution}{exo:g8:lenses-images:6}
\omterm{def:g8:lenses-images:families}{透镜}：玻璃物镜——眼睛中位于瞳孔后方的
晶状体。屏幕：感光元件——视网膜。调焦：相机
通过移动镜头；眼睛则改变晶状体形状，看近物时变厚，
看远物时变薄。
\end{solution}

\begin{solution}{exo:g8:lenses-images:7}
近视眼聚焦过强，像在视网膜之前形成——
此时用\omterm{def:g8:lenses-images:families}{凹透镜}稍稍减弱会聚作用即可。年老的远视眼
看近处时聚焦不足，像落在视网膜之后；
应配戴\omterm{def:g8:lenses-images:families}{凸透镜}，补足所缺的会聚能力。
\end{solution}

\begin{solution}{exo:g8:lenses-images:8}
邮票必须放在比\omterm{def:g8:lenses-images:focus}{焦距}更靠近\omterm{def:g8:lenses-images:families}{透镜}的位置。
此时眼睛看到的是正立、放大的虚像——\omterm{def:g6:light-propagation:ray}{光线}的方向
仿佛来自玻璃后方一枚更大的邮票；
这种像不能落在任何屏幕上。
\end{solution}

\begin{solution}{exo:g8:lenses-images:9}
物体靠近时，像向后退——因此相机
要把镜头\emph{远离}感光元件，以追随后退的像。
眼睛则通过使晶状体变厚，
弯曲得更厉害，把像拉回固定的视网膜上。
\end{solution}

\begin{solution}{exo:g8:lenses-images:10}
只有在\omterm{def:g8:lenses-images:focus}{焦距}处，\omterm{def:g8:lenses-images:families}{透镜}才能把近似平行的太阳光
会聚成最小、最炽烈的光斑。
无论更近或更远，\omterm{def:g6:light-propagation:ray}{光线}都会散布在一块过于柔和、
无法引燃的区域上——火种恰好在距离 $f$ 处。
\end{solution}

\begin{solution}{exo:g8:lenses-images:11}
老花镜是\omterm{def:g8:lenses-images:families}{凸透镜}：它们能把阳光会聚到一个可以点燃的\omterm{def:g8:lenses-images:focus}{焦点}上，
还能放大近处的小字——这两种本领都只属于\omterm{def:g8:lenses-images:families}{凸透镜}。
而近视的孙子所戴的\omterm{def:g8:lenses-images:families}{凹透镜}则使每束\omterm{def:g6:light-propagation:ray}{光线}都向外散开：
没有\omterm{def:g8:lenses-images:focus}{焦点}，也没有能点燃的光斑，透过它们看到的
字不但不会变大，反而缩小了。
\end{solution}

\begin{solution}{exo:g8:lenses-images:12}
望远镜第一级：大\omterm{def:g8:lenses-images:families}{透镜}接收\omterm{def:g4:solar-system:star}{恒星}近似平行的\omterm{def:g6:light-propagation:ray}{光线}，
使其会聚，在\omterm{def:g8:lenses-images:focus}{焦点}附近形成小的\omterm{prop:g8:lenses-images:realimage}{实像}；第二级再放大该像。
微弱的星光十分珍贵——前端\omterm{def:g8:lenses-images:families}{透镜}的\emph{口径}
决定每颗\omterm{def:g4:solar-system:star}{恒星}能收集多少光，口径几乎决定一切。
显微镜第一级：弯曲强烈、$f$ 较小的\omterm{def:g8:lenses-images:families}{透镜}，
紧贴液滴悬放，形成放大很多的\omterm{prop:g8:lenses-images:realimage}{实像}；
第二级再次放大。它的物体明亮且距离很近，
第一级所需的不是口径，
而是弯曲的本领。
同一结构，需求却截然相反。
\end{solution}

\begin{solution}{pb:g8:lenses-images:1}
\textbf{1.} 触碰：中间鼓起的为\omterm{def:g8:lenses-images:families}{凸透镜}；中间凹陷的，
为\omterm{def:g8:lenses-images:families}{凹透镜}。窗口测试：\omterm{def:g8:lenses-images:families}{凸透镜}在某一距离处，
能在墙上投出倒立的窗口像；
\omterm{def:g8:lenses-images:families}{凹透镜}则永远做不到。
\textbf{2.} $f = \qty{25}{cm}$；\omterm{def:g8:lenses-images:families}{凸透镜}。
\textbf{3.} \omterm{def:g8:lenses-images:families}{凹透镜}——在任何距离都不成\omterm{prop:g8:lenses-images:realimage}{实像}，
透过它看到的缩小景象已足以下结论。
\textbf{4.} C：\omterm{def:g8:lenses-images:focus}{焦距}较短（\qty{10}{cm}）者，
说明\omterm{def:g8:lenses-images:families}{透镜}更厚，弯曲也更厉害。
\textbf{5.} 远处模糊、近处清楚：近视，
即眼睛弯曲过强，应配戴\omterm{def:g8:lenses-images:families}{凹透镜}类
的眼镜。
\textbf{6.} 在一臂距离处阅读困难：年老的远视眼，
聚焦较弱，应配戴\omterm{def:g8:lenses-images:families}{凸透镜}阅读镜。
\textbf{7.} 阅读镜所增加的会聚能力，是为近处书页设计的；
若对着远山观看，山上的\omterm{def:g6:light-propagation:ray}{光线}近似平行，无需帮助，
所增加的能力便会矫枉过正——像会在视网膜之前成形，
山峰因此模糊。
\textbf{8.} 为视力良好的眼睛加上矫正镜，会使其失衡：
试戴的\omterm{def:g8:lenses-images:families}{透镜}使孩子的眼睛须与过度弯曲的\omterm{def:g6:light-propagation:ray}{光线}抗衡，
像落在视网膜之外，视力表上的字似在游动，
正是健康眼睛的抗议。
\textbf{9.} 前端\omterm{def:g8:lenses-images:families}{透镜}把\omterm{def:g4:solar-system:star}{恒星}的平行\omterm{def:g6:light-propagation:ray}{光线}会聚，
在其\omterm{def:g8:lenses-images:focus}{焦点}附近形成一个小的\omterm{prop:g8:lenses-images:realimage}{实像}；目镜是一个放大镜，
放在该\omterm{prop:g8:lenses-images:realimage}{实像}的\omterm{def:g8:lenses-images:focus}{焦距}之内，
把\omterm{prop:g8:lenses-images:realimage}{实像}放大后送入观测者的眼睛。
\textbf{10.} 每颗\omterm{def:g4:solar-system:star}{恒星}每平方\omterm{def:g2:measuring-length:units}{厘米}仅送来一丝微光；
碟子般大的\omterm{def:g8:lenses-images:families}{透镜}，可一次收集数百平方\omterm{def:g2:measuring-length:units}{厘米}，
将整整一只餐盘大小的微弱\omterm{def:g6:light-propagation:ray}{光线}汇聚成一个明亮的光点。
口径正是天文学家所渴求的。
\textbf{11.} 硬币大小的放大镜，即使弯曲再厉害，
也只能收集一枚硬币大小的星光，而所需的是碟子大小；
像或许清晰，却光量不足——
暗星因此消失。它的收获不过是目镜本已能提供的，
结论：不能替代，应另订合适的\omterm{def:g8:lenses-images:families}{透镜}。
\textbf{12.} 例如：“货架上的每件东西都在做两件事：
\emph{收集}光，再把它\emph{送}到眼睛所需之处——
而文明的整套光学工具，从眼镜到天文台，
都只是这两个动作的不同规模，
尺寸各异而已。”
\end{solution}
